\chapter{Contribution Guide}
\label{ch:contribution-guide}

\section{Basis and honesty note}

The repository contains \textbf{no} \texttt{CONTRIBUTING.md}, no
\texttt{CODE\_OF\_CONDUCT.md}, no issue or pull-request templates and no
\texttt{GOVERNANCE} file. There is therefore no author-stated contribution
process to document. What follows is derived from the repository's structure,
scripts and workflows: it is a factual description of how work can be
contributed, labelled INFERRED, followed by recommended additions labelled
RECOMMENDED. It is not a substitute for a real contribution policy.

\section{Prerequisites (INFERRED)}

\begin{itemize}
  \item Node.js 18 or newer (20 recommended; CI uses 20).
  \item pnpm 9 or newer.
  \item Git.
  \item Optional: the VS Code Extension Development Host for working on the
        extension.
\end{itemize}

\section{Getting started (INFERRED)}

\begin{lstlisting}[language=bash]
git clone <repository>
cd expo-snips
pnpm install
pnpm build            # build all packages
pnpm validate         # check every recipe against the schema
\end{lstlisting}

To run the registry locally:

\begin{lstlisting}[language=bash]
pnpm serve            # builds the registry and serves it on PORT (default 3000)
\end{lstlisting}

To work on the extension:

\begin{lstlisting}[language=bash]
pnpm --filter expo-snips build       # bundle to out/extension.js
pnpm --filter expo-snips watch       # rebuild on change
pnpm --filter expo-snips typecheck   # no-emit type check
\end{lstlisting}

Press F5 with the extension package open to launch an Extension Development
Host.

\section{Contribution paths}

There are two natural kinds of contribution, and they differ substantially.

\subsection*{1. Adding or changing a screen recipe (most common)}

\begin{enumerate}
  \item Create or edit a JSON document under
        \texttt{packages/registry/data/<category>/<name>.json}. The file path must
        match the recipe id, and the document must satisfy
        \texttt{packages/registry/schema.json}.
  \item Add the templates under
        \texttt{packages/registry/templates/<category>/<name>/}, using the dotted
        naming convention (see Chapter~\ref{ch:development-guidelines}).
  \item Declare every third-party package the templates import in the recipe's
        \texttt{dependencies}.
  \item Run \texttt{pnpm validate} to confirm the recipe is well-formed and all
        referenced templates exist.
  \item Run \texttt{pnpm build:registry} (or \texttt{pnpm serve}) to confirm the
        recipe builds and to update the committed \texttt{registry/} output.
\end{enumerate}

The validator gives fast, specific feedback (id format, version format, enum
values, template existence, id/path agreement), so this workflow is contributor
friendly.

\subsection*{2. Changing engine or client code}

\begin{enumerate}
  \item Make the change in the relevant package (\texttt{core} for the shared
        engine, \texttt{cli} or \texttt{vscode-extension} for a client).
  \item Run \texttt{pnpm build} to type-check and compile.
  \item Manually verify generation, since \textbf{no automated tests exist}
        (Chapter~\ref{ch:testing}).
  \item If the behaviour is user-visible, note that the same behaviour may need
        to be applied in both clients, because orchestration is duplicated
        (Chapter~\ref{ch:technical-debt}, TD-08).
\end{enumerate}

\section{Change and release process (INFERRED)}

\begin{itemize}
  \item Commits are made to \texttt{master}. Existing history shows small,
        direct commits (ten commits at the time of analysis) rather than a
        feature-branch/PR flow.
  \item The registry is redeployed automatically on every push to
        \texttt{main} or \texttt{master}.
  \item The npm packages are released by pushing a tag matching \texttt{v*}.
  \item Version numbers are edited by hand in the relevant
        \texttt{package.json}.
  \item The extension is packaged by hand
        (\texttt{pnpm --filter expo-snips package}).
\end{itemize}

\textbf{Consequence:} contributing directly to the default branch is effectively
the current practice, and there is no automated gate before a change reaches the
deployed registry. A contributor should therefore run \texttt{pnpm validate} and
\texttt{pnpm build} locally before pushing, because CI will not.

\section{What to include in a contribution}

\begin{longtable}{p{4.2cm}p{10cm}}
\caption{Contribution expectations (INFERRED)}\\
\toprule
\textbf{Change type} & \textbf{Expected accompanying work} \\
\midrule

New recipe & JSON document, templates, dependency declarations, validator run,
registry rebuild. \\
Recipe change & Version bump of the recipe's own \texttt{version} if the change is
behavioural; note whether existing projects need to regenerate. \\
Core engine change & Manual verification across several recipes and both clients;
consider a test once tests exist. \\
Client change & Verify the same behaviour is not required in the other client. \\
Extension release & Update \texttt{CHANGELOG.md}; the top of the file establishes
Keep a Changelog and SemVer. \\
Breaking change & No deprecation policy exists; document the impact explicitly
(see below). \\
\bottomrule
\end{longtable}

\section{Missing policy areas}

\textbf{Status:} RECOMMENDED. The following should be created before the project
invites external contributors:

\begin{enumerate}
  \item A \texttt{CONTRIBUTING.md} capturing the above, plus branch and PR
        conventions.
  \item A \texttt{CODE\_OF\_CONDUCT.md}.
  \item Issue and pull-request templates.
  \item A documented deprecation and recipe-migration policy, since recipes are
        versioned independently and users may have generated screens from old
        versions.
  \item A root \texttt{README} that orients a newcomer (currently absent).
  \item A licence decision for the extension: it currently ships under a
        ``Private License'', which will deter external contribution. This is a
        deliberate-looking author choice, but it should be stated plainly for
        contributors.
\end{enumerate}

\section{Contribution-friendly features already present}

\begin{itemize}
  \item The recipe validator provides immediate, actionable feedback.
  \item Recipes are independent files, so two contributors rarely conflict.
  \item The offline build means a contributor can fully test generation without
        any network access.
  \item The extension CHANGELOG already follows a recognised convention.
  \item The extension README documents the source layout and a build/debug
        workflow.
\end{itemize}

\textbf{Status:} IMPLEMENTED.

\section{Summary}

\textbf{Status:} INFERRED (all process descriptions) / RECOMMENDED (formal
contribution policy). The technical workflow for adding a recipe is genuinely
good; the social and policy scaffolding for outside contributors does not yet
exist.
